% part: intuitionistic-logic
% chapter: introduction
% section: natural-deduction

\documentclass[../../../include/open-logic-section]{subfiles}

\begin{document}

\olfileid[pt]{int}{int}{ntd}

\olsection{Dedução natural}

A dedução natural sem as regras $\FalseCl$ é um sistema padrão de
!!{derivation} para a lógica intuicionista. Repetimos aqui as regras e
indicamos sua motivação mediante a interpretação BHK. Em cada caso,
podemos pensar em uma regra que nos permite concluir que, se as
premissas possuem construções, a conclusão também possui.

Como as !!{derivation}s de dedução natural têm hipóteses não
descartadas, devemos entender tal !!a{derivation}, digamos, de $!A$ a
partir das hipóteses~$\Gamma$ !!{undischarged}, como uma função que
transforma construções de todos os $!B \in \Gamma$ em uma construção
de~$!A$. Se existe !!a{derivation} de~$!A$ sem hipóteses
!!{undischarged}, então existe uma construção de~$!A$ no sentido da
interpretação BHK. Para fins desta discussão, porém, omitiremos
$\Gamma$ quando ela não for necessária.

Uma hipótese~$!A$, por si só, é !!a{derivation} de~$!A$ a partir da
hipótese~$!A$ !!{undischarged}. Isso está de acordo com a interpretação
BHK: a função identidade sobre construções transforma qualquer
construção de~$!A$ em uma construção de~$!A$.

\subsection{Conjunção}

\begin{defish}
\AxiomC{$!A$}
\AxiomC{$!B$}
\RightLabel{\Intro{\land}}
\BinaryInfC{$!A \land !B$}
\DisplayProof
\hfill
\begin{tabular}{l}
\AxiomC{$!A \land !B$}
\RightLabel{\Elim{\land}}
\UnaryInfC{$!A$}
\DisplayProof
\\[3ex]
\AxiomC{$!A \land !B$}
\RightLabel{\Elim{\land}}
\UnaryInfC{$!B$}
\DisplayProof
\end{tabular}
\end{defish}

Suponha que tenhamos construções $N_1$, $N_2$ de $!A_1$ e $!A_2$,
respectivamente. Então também temos uma construção de
$!A_1 \land !A_2$, a saber, o par $\tuple{N_1, N_2}$.

Uma construção de $!A_1 \land !A_1$ segundo a interpretação BHK é um
par $\tuple{N_1, N_2}$. Suponha, portanto, que tenhamos tal par. Então
também temos uma construção de cada membro da conjunção: $N_1$ é uma construção
de~$!A_1$ e $N_2$ é uma construção de $!A_2$.

\subsection{Condicional}

\begin{defish}
  \AxiomC{$\Discharge{!A}{u}$}
  \DeduceC{$!B$}
  \DischargeRule{$\Intro{\lif}$}{u}
  \UnaryInfC{$!A \lif !B$}
  \DisplayProof
\hfill
  \AxiomC{$!A \lif !B$}
  \AxiomC{$!A$}
  \RightLabel{$\Elim{\lif}$}
  \BinaryInfC{$!B$}
  \DisplayProof
\end{defish}

Se temos !!a{derivation} de~$!B$ a partir da hipótese~$!A$
!!{undischarged}, então existe uma função~$f$ que transforma construções
de~$!A$ em construções de~$!B$. Essa mesma função é uma construção de
$!A \lif !B$. Assim, se a premissa de $\Intro\lif$ possui uma construção
condicionada a uma construção de~$!A$, a conclusão $!A \lif !B$ possui
uma construção.

Por outro lado, suponha que existam construções~$N$ de $!A$ e $f$ de
$!A \lif !B$. Uma construção de $!A \lif !B$ é uma função que
transforma construções de $!A$ em construções de~$!B$. Portanto,
$f(N)$ é uma construção de~$!B$, isto é, a conclusão de $\Elim\lif$
possui uma construção.

\subsection{Disjunção}

\begin{defish}
\begin{tabular}{l}
\AxiomC{$!A$}
\RightLabel{\Intro{\lor}}
\UnaryInfC{$!A \lor !B$}
\DisplayProof
\\[3ex]
\AxiomC{$!B$}
\RightLabel{\Intro{\lor}}
\UnaryInfC{$!A \lor !B$}
\DisplayProof
\end{tabular}
\hfill
\AxiomC{$!A \lor !B$}
\AxiomC{$\Discharge{!A}{n}$}
\DeduceC{$!C$}
\AxiomC{$\Discharge{!B}{n}$}
\DeduceC{$!C$}
\DischargeRule{\Elim{\lor}}{n}
\TrinaryInfC{$!C$}
\DisplayProof
\end{defish}

Se temos uma construção~$N_i$ de~$!A_i$, podemos transformá-la em uma
construção~$\tuple{i, N_i}$ de~$!A_1 \lor !A_2$. Por outro lado,
suponha que tenhamos uma construção de~$!A_1 \lor !A_2$, isto é, um
par $\tuple{i, N_i}$ em que $N_i$ é uma construção de~$!A_i$, e também
funções $f_1$, $f_2$ que transformam, respectivamente, construções de
$!A_1$, $!A_2$ em construções de~$!C$. Então $f_i(N_i)$ é uma
construção de~$!C$, a conclusão de~$\Elim\lor$.

\subsection{Absurdo}

\begin{defish}
  \AxiomC{$\lfalse$}
  \RightLabel{\FalseInt}
  \UnaryInfC{$!A$}
  \DisplayProof
\end{defish}

Se temos !!a{derivation} de~$\lfalse$ a partir das hipóteses
$!B_1$, \dots, $!B_n$ !!{undischarged}, então existe uma
função $f(M_1, \dots, M_n)$ que transforma construções de $!B_1$,
\dots, $!B_n$ em uma construção de~$\lfalse$. Como $\lfalse$ não possui
construção, também não podem existir construções de todos os $!B_1$,
\dots, $!B_n$. Logo, $f$ também possui a propriedade de que \emph{se}
$M_1$, \dots, $M_n$ são construções de $!B_1$, \dots, $!B_n$,
respectivamente, \emph{então} $f(M_1, \dots, M_n)$ é uma construção
de~$!A$.

\subsection{Regras para $\lnot$}

Como $\lnot !A$ é definida como $!A \lif \lfalse$, estritamente falando
não precisamos de regras para~$\lnot$. Mas, se precisássemos, elas
teriam a seguinte forma:

\begin{defish}
\AxiomC{$\Discharge{!A}{n}$}
\noLine
\DeduceC{$\lfalse$}
\DischargeRule{\Intro{\lnot}}{n}
\UnaryInfC{$\lnot !A$}
\DisplayProof
\hfill
\AxiomC{$\lnot !A$}
\AxiomC{$!A$}
\RightLabel{\Elim{\lnot}}
\BinaryInfC{$\lfalse$}
\DisplayProof
\end{defish}


\subsection{Exemplos de \usetoken{P}{derivation}}

\begin{enumerate}
\item $\Proves !A \lif (\lnot !A \lif \lfalse)$, isto é, $\Proves !A
  \lif ((!A \lif \lfalse) \lif \lfalse)$
  \begin{prooftree}
    \AxiomC{$\Discharge{!A}{2}$}
    \AxiomC{$\Discharge{!A \lif \lfalse}{1}$}
    \RightLabel{\Elim\lif}
    \BinaryInfC{$\lfalse$}
    \DischargeRule{\Intro\lif}{1}
    \UnaryInfC{$(!A \lif \lfalse)\lif \lfalse$}
    \DischargeRule{\Intro\lif}{2}
    \UnaryInfC{$!A \lif (!A \lif \lfalse) \lif \lfalse$}
  \end{prooftree}

\item $\Proves ((!A \land !B) \lif !C) \lif (!A \lif (!B \lif !C))$
  \begin{prooftree}
    \AxiomC{$\Discharge{(!A \land !B) \lif !C}{3}$}
    \AxiomC{$\Discharge{!A}{2}$}
    \AxiomC{$\Discharge{!B}{1}$}
    \RightLabel{\Intro\land}
    \BinaryInfC{$!A \land !B$}
    \RightLabel{\Elim\lif}
    \BinaryInfC{$!C$}
    \DischargeRule{\Intro\lif}{1}
    \UnaryInfC{$!B \lif !C$}
    \DischargeRule{\Intro\lif}{2}
    \UnaryInfC{$!A \lif (!B \lif !C)$}
    \DischargeRule{\Intro\lif}{3}
    \UnaryInfC{$((!A \land !B) \lif !C) \lif (!A \lif (!B \lif !C))$}
  \end{prooftree}

\item $\Proves \lnot(!A \land \lnot !A)$, isto é, $\Proves (!A \land (!A
  \lif \lfalse))\lif \lfalse$
  \begin{prooftree}
    \AxiomC{$\Discharge{!A \land (!A \lif \lfalse)}{1}$}
    \RightLabel{\Elim\land}
    \UnaryInfC{$!A \lif \lfalse$}
    \AxiomC{$\Discharge{!A \land (!A \lif \lfalse)}{1}$}
    \RightLabel{\Elim\land}
    \UnaryInfC{$!A$}
    \RightLabel{\Elim\lif}
    \BinaryInfC{$\lfalse$}
    \DischargeRule{\Intro\lif}{1}
    \UnaryInfC{$(!A \land (!A \lif \lfalse))\lif \lfalse$}
  \end{prooftree}

\item $\Proves \lnot\lnot(!A \lor \lnot !A)$, isto é, $\Proves ((!A \lor
  (!A \lif \lfalse))\lif \lfalse)\lif \lfalse$
  \begin{prooftree}\hskip-3em
    \AxiomC{$\Discharge{(!A \lor (!A \lif \lfalse))\lif \lfalse}{2}$}
    \AxiomC{$\Discharge{(!A \lor (!A \lif \lfalse))\lif \lfalse}{2}$}
    \AxiomC{$\Discharge{!A}{1}$}
    \RightLabel{\Intro\lor}
    \UnaryInfC{$!A \lor (!A \lif \lfalse)$}
    \RightLabel{\Elim\lif}
    \BinaryInfC{$\lfalse$}
    \DischargeRule{\Intro\lif}{1}
    \UnaryInfC{$!A \lif \lfalse$}
    \RightLabel{\Intro\lor}
    \UnaryInfC{$!A \lor (!A \lif \lfalse)$}
    \RightLabel{\Elim\lif}
    \insertBetweenHyps{\hskip -4em}
    \BinaryInfC{$\lfalse$}
    \DischargeRule{\Intro\lif}{2}
    \UnaryInfC{$((!A \lor (!A \lif \lfalse))\lif \lfalse)\lif \lfalse$}
  \end{prooftree}
\end{enumerate}

\begin{prop}
  Se $\Gamma \Proves !A$ na dedução natural para a lógica
  intuicionista, então $\Gamma \Proves !A$ na lógica clássica. Em
  particular, se $!A$ é um teorema intuicionista, também é um teorema
  clássico.
\end{prop}

\begin{proof}
  Toda regra de dedução natural também é uma regra da dedução natural
  clássica; portanto, toda !!{derivation} na lógica intuicionista também
  é !!a{derivation} na lógica clássica.
\end{proof}

\begin{prob}
  Forneça !!{derivation}s na lógica intuicionista das seguintes !!{formula}s:
  \begin{enumerate}
    \item $(\lnot !A \lor !B) \lif (!A \lif !B)$
    \item $\lnot\lnot\lnot !A \lif \lnot !A$
    \item $\lnot\lnot (!A \land !B) \liff (\lnot\lnot !A\land \lnot \lnot !B)$
    \item $\lnot (!A \lor !B) \liff (\lnot !A \land \lnot !B)$
    \item $(\lnot !A \lor \lnot !B) \lif \lnot (!A \land !B)$
    \item $\lnot\lnot ( !A \land !B) \lif  (\lnot\lnot !A \lor
    \lnot\lnot !B)$
  \end{enumerate}
\end{prob}

\end{document}
